% ==============================================================================
% Capítulo 2: Fundamentação teórica
% ==============================================================================
\chapter{Fundamentação teórica}
\label{cap:fundamentacao}

\begin{instrucao}
\orientacao{O que vai na Fundamentação Teórica}
Aqui entram os conceitos de que quem lê o trabalho precisa para entender o restante do
trabalho, e somente eles. A Fundamentação não é um livro-texto sobre o
assunto. Antes de incluir um conceito, pergunte se ele será usado nos
capítulos seguintes. Como referência, de dez a quinze páginas costumam ser
suficientes.

Escreva com suas palavras e cite as fontes (livros, artigos, teses). A citação
direta (cópia literal) deve ser exceção, reservada a definições ou trechos
cuja redação exata importa. Uma sequência de citações diretas não constitui
fundamentação. Copiar texto sem citar é plágio e, pelas regras do TCC, leva à
reprovação com nota zero.

Abra o capítulo com um parágrafo curto que diga o que será apresentado e em
que ordem. Evite seções com uma única subseção e títulos sem texto entre eles.
\end{instrucao}

Este capítulo apresenta os conceitos que fundamentam o trabalho. A
\cref{sec:fundamentacao:mkp} define formalmente o \acs{mkp}. A
\cref{sec:fundamentacao:pli} trata da resolução exata por \acl{pli} e de suas
limitações. A \cref{sec:fundamentacao:ils} apresenta a metaheurística de
\acl{ils}, base da heurística proposta.

% ------------------------------------------------------------------------------
\section{O problema da mochila multidimensional}
\label{sec:fundamentacao:mkp}

Considere um conjunto de $n$ itens e $m$ recursos. Cada item $j \in \{1, \ldots, n\}$
tem um valor $p_j > 0$ e consome uma quantidade $w_{ij} \geq 0$ de cada recurso
$i \in \{1, \ldots, m\}$, cuja capacidade total é $c_i > 0$. Seja $x_j$ uma
variável binária que vale 1 se o item $j$ for escolhido e 0 caso contrário. O
\acs{mkp} pode ser formulado como o modelo de \acs{pli} dado pelas
\crefrange{eq:mkp:objetivo}{eq:mkp:dominio} \cite{kellerer2004knapsack}.
\begin{align}
  \max \quad & \sum_{j=1}^{n} p_j \, x_j \label{eq:mkp:objetivo} \\
  \text{sujeito a} \quad & \sum_{j=1}^{n} w_{ij} \, x_j \leq c_i,
      && i = 1, \ldots, m, \label{eq:mkp:capacidade} \\
  & x_j \in \{0, 1\}, && j = 1, \ldots, n. \label{eq:mkp:dominio}
\end{align}

A \cref{eq:mkp:objetivo} maximiza o valor total dos itens escolhidos. As
restrições~\eqref{eq:mkp:capacidade} garantem que a capacidade de cada recurso
seja respeitada e as restrições~\eqref{eq:mkp:dominio} definem o domínio das
variáveis. Quando $m = 1$, o modelo se reduz ao problema da mochila clássico,
que pode ser resolvido em tempo pseudopolinomial por programação dinâmica. Para
$m \geq 2$, o problema é fortemente NP-difícil \cite{kellerer2004knapsack}.

\begin{instrucao}
\orientacao{Equações}
Use \verb|equation| para uma equação e \verb|align| para várias alinhadas
(o \verb|&| marca o ponto de alinhamento). Toda equação citada no texto
precisa de \verb|\label|. Para não numerar uma linha, use \verb|\notag|.
\begin{exemplo}
A área do círculo é dada pela
\cref{eq:area}.
\begin{equation}
  A = \pi r^2,
  \label{eq:area}
\end{equation}
em que $r$ é o raio.
\end{exemplo}
Algumas regras:
\begin{itemize}
  \item a equação faz parte da frase, por isso termina com vírgula ou ponto
        conforme o texto que a segue;
  \item defina cada símbolo na primeira vez que ele aparece;
  \item use a vírgula como separador decimal, também nas fórmulas: \verb|$3,14$|
        produz $3,14$ (o pacote \verb|icomma| evita o espaço depois da vírgula);
  \item para várias equações use \verb|\crefrange{eq:a}{eq:c}|, que gera
        ``Equações (1) a (3)''.
\end{itemize}
\end{instrucao}

% ------------------------------------------------------------------------------
\section{Programação linear inteira}
\label{sec:fundamentacao:pli}

Modelos de \acs{pli} são resolvidos, em geral, pelo método \emph{branch-and-bound}
e suas variantes, que dividem o problema em subproblemas e usam a relaxação
linear de cada um para descartar regiões do espaço de busca que não podem
conter a solução ótima \cite{wolsey2020integer}. A relaxação linear é obtida
substituindo a restrição~\eqref{eq:mkp:dominio} por $0 \leq x_j \leq 1$, o que
transforma o modelo em um problema de programação linear, resolvível em tempo
polinomial.

A eficiência desse processo depende da qualidade dos limitantes obtidos. No
\acs{mkp}, a diferença entre o valor da relaxação linear e o da solução ótima
tende a diminuir com o aumento do número de itens, mas o número de subproblemas
explorados cresce rapidamente com o número de restrições
\cite{puchinger2010multidimensional}. Solucionadores modernos, como o HiGHS
\cite{huangfu2018parallelizing}, combinam esse método com planos de corte e
heurísticas primais, o que os torna competitivos em instâncias pequenas e
médias.

\begin{instrucao}
\orientacao{Citação direta}
A citação direta reproduz literalmente o texto da fonte e sempre informa a
página. Com até três linhas, fica no próprio parágrafo, entre aspas duplas.
\begin{exemplo}
A portaria define alucinação como
``informação aparentemente plausível,
porém incorreta, imprecisa ou
inventada'' \cite[p.~1]{ufes2026portaria}.
\end{exemplo}
Com mais de três linhas, use o ambiente \verb|citacao|, que aplica recuo de
4~cm, fonte menor e espaçamento simples, sem aspas.
\begin{exemplo}
\begin{citacao}
O usuário é responsável pela precisão,
pela originalidade e pela adequação dos
conteúdos gerados ou modificados por
sistema de inteligência artificial que
utilizar, devendo verificá-los antes de
sua utilização, considerada a
possibilidade de alucinação
\cite[p.~5]{ufes2026portaria}.
\end{citacao}
\end{exemplo}
Indique supressões no meio do trecho com \verb|[...]| e destaques seus com a
expressão ``grifo nosso'' após a página.
\end{instrucao}

% ------------------------------------------------------------------------------
\section{Metaheurísticas e busca local iterada}
\label{sec:fundamentacao:ils}

Metaheurísticas são estratégias gerais que guiam heurísticas subordinadas na
exploração do espaço de soluções, com o objetivo de escapar de ótimos locais.
Uma das mais simples é a \ac{ils}, que alterna duas fases: uma busca local, que
melhora a solução corrente até atingir um ótimo local, e uma perturbação, que
modifica esse ótimo local para gerar um novo ponto de partida
\cite{lourenco2019iterated}. Se a busca local a partir do ponto perturbado
encontrar uma solução melhor, ela passa a ser a solução corrente. Caso
contrário, um critério de aceitação decide se a busca continua a partir da
nova solução ou retorna à anterior.

A \cref{fig:ils} ilustra esse comportamento em um problema de maximização. A
partir de uma solução inicial $s_0$, a busca local sobe até o ótimo local
$s^*$. A perturbação leva a busca para $s'$, fora da região de atração de
$s^*$, e uma nova busca local alcança $s^{*\prime}$, de valor maior. A
intensidade da perturbação é o parâmetro central do método. Perturbações
fracas fazem a busca local retornar ao mesmo ótimo e perturbações fortes
equivalem a reiniciar a busca de um ponto aleatório
\cite{lourenco2019iterated}.

\begin{figure}[htbp]
  \centering
  \caption{Funcionamento da \acs{ils} em um problema de maximização}
  \label{fig:ils}
  \begin{tikzpicture}
    \begin{axis}[
        width=0.85\textwidth, height=6cm,
        domain=0:10, samples=200,
        axis lines=left, xtick=\empty, ytick=\empty,
        xlabel={Espaço de soluções}, ylabel={Valor da solução},
        ymin=1.8, ymax=6.2, clip=false,
      ]
      \addplot[thick, black] {3 + sin(deg(x)) + 0.5*sin(deg(2.7*x)) + 0.15*x};
      \node[circle, fill=black, inner sep=1.6pt, label=below:{$s_0$}] (s0) at (axis cs:4.6,2.62) {};
      \node[circle, fill=black, inner sep=1.6pt, label=above:{$s^*$}] (s1) at (axis cs:5.49,3.50) {};
      \node[circle, draw=black, fill=white, inner sep=1.6pt, label=below right:{$s'$}] (s2) at (axis cs:6.5,3.71) {};
      \node[circle, fill=black, inner sep=1.6pt, label=above:{$s^{*\prime}$}] (s3) at (axis cs:7.66,5.61) {};
      \draw[-{Stealth}, thick, gray] (s0) to[bend right=25] node[right, font=\footnotesize, black] {busca local} (s1);
      \draw[-{Stealth}, thick, dashed] (s1) to[bend left=40] node[above, font=\footnotesize] {perturbação} (s2);
      \draw[-{Stealth}, thick, gray] (s2) to[bend right=25] node[right, font=\footnotesize, black] {busca local} (s3);
    \end{axis}
  \end{tikzpicture}
  \fonteautor
\end{figure}

\begin{instrucao}
\orientacao{Figuras: inclusão, citação e qualidade}
Pela NBR 14724, a identificação da figura (a palavra ``Figura'', o número e
o título) fica \textbf{acima} dela e a fonte fica \textbf{abaixo}, mesmo quando a figura foi
feita por você. Por isso, no LaTeX, o \verb|\caption| vem antes do conteúdo e
o \verb|\fonte| depois. Para material próprio, use \verb|\fonteautor|, que
escreve ``Fonte: Elaborado pelo próprio autor''. Para figura adaptada de outra
obra, escreva \verb|\fonte{Adaptado de \textcite{chave}.}|.
\begin{exemplocodigo}
\begin{figure}[htbp]
  \centering
  \caption{Arquitetura do sistema}
  \label{fig:arquitetura}
  \includegraphics[width=0.8\textwidth]{figuras/arquitetura.pdf}
  \fonteautor
\end{figure}
\end{exemplocodigo}

\textbf{Toda figura, tabela, quadro, código ou algoritmo incluído no trabalho
deve ser citado no texto pelo menos uma vez}, de preferência antes de aparecer.
Ao citar, escreva o nome com inicial maiúscula seguido do número, como em
``a Figura 3 mostra'', e não ``a figura 3''. O mesmo vale para Tabela, Quadro,
Código, Algoritmo, Equação e Seção. Os comandos \verb|\cref| e \verb|\Cref| já
escrevem o nome dessa forma. Não basta citar. Explique no texto o que deve ser observado na
figura. Gráficos e diagramas mais complexos pedem um ou dois parágrafos de
explicação, como foi feito acima para a \cref{fig:ils}.

\textbf{Como obter figuras de boa qualidade}
\begin{itemize}
  \item Prefira figuras \textbf{vetoriais} (PDF, EPS ou desenhadas em TikZ),
        que não perdem nitidez com o zoom nem na impressão. Esta figura foi
        desenhada em TikZ diretamente no texto.
  \item Gráficos feitos em Python devem ser salvos em PDF:
        \verb|plt.savefig("grafico.pdf", bbox_inches="tight")|. Guarde também o
        script que gerou o gráfico, para poder refazê-lo.
  \item Diagramas feitos no draw.io (diagrams.net), no Inkscape ou em
        ferramentas UML devem ser exportados em PDF ou SVG (convertido para
        PDF), nunca capturados com \emph{print screen}.
  \item Capturas de tela de sistemas devem ser salvas em PNG, na resolução
        original e com o zoom do navegador em 100\%. Evite JPEG em diagramas e
        capturas de tela, porque a compressão borra textos e linhas.
  \item Ajuste apenas a largura (\verb|width=|). Alterar largura e altura ao
        mesmo tempo distorce a imagem.
  \item O texto dentro da figura deve ter tamanho próximo ao das legendas.
        Uma figura com letras minúsculas foi feita grande demais e reduzida.
  \item Use cores que continuem distinguíveis em escala de cinza e para
        pessoas com daltonismo. Variar também o tipo de linha e de marcador
        ajuda.
  \item Mantenha o mesmo estilo (fonte, espessura de linhas, cores) em todas
        as figuras do trabalho e o mesmo alinhamento (centralizado), como pede
        a NBR 14724:2024.
\end{itemize}
\end{instrucao}

A eficácia da \acs{ils} depende, portanto, da busca local, da perturbação e
do critério de aceitação, componentes que precisam ser ajustados a cada
problema. O \cref{cap:metodologia} descreve como cada um deles foi definido
para o \acs{mkp}.
